\documentclass[12pt,a4paper]{article}

\usepackage[T1]{fontenc}
\usepackage[polish]{babel}
\usepackage[a4paper,margin=2.5cm]{geometry}
\usepackage{graphicx}
\usepackage{hyperref}
\usepackage{amssymb}
\usepackage{amsmath}
\usepackage{booktabs}
\usepackage{enumitem}
\usepackage{listings}
\usepackage{xcolor}

% --- listings: prosty styl dla fragmentów kodu / poleceń ---
\lstset{
    basicstyle=\ttfamily\small,
    breaklines=true,
    frame=single,
    columns=fullflexible,
    keepspaces=true,
    showstringspaces=false,
}

% --- pomocnicze makro na znaczniki TODO (widoczne na czerwono) ---
\newcommand{\todo}[1]{\textcolor{red}{\textbf{[TODO: #1]}}}

\newcommand{\coursename}{Wstęp do algorytmów ewolucyjnych}
\newcommand{\projecttitle}{Harmonogramowanie zadań}
\newcommand{\doctype}{Dokumentacja końcowa}
\newcommand{\authorone}{Wojciech Jasiński}
\newcommand{\authortwo}{Piotr Iśtok}
\newcommand{\projectdate}{Warszawa, \the\year}

\begin{document}
\sloppy

\begin{titlepage}
    \centering

    \includegraphics[width=0.8\textwidth]{politechnika.png}\par
    \vspace{2.5cm}

    {\Large \coursename\par}
    \vspace{1.2cm}

    {\Huge \bfseries \projecttitle \par}
    \vspace{1cm}

    {\Large \itshape \doctype \par}
    \vspace{2cm}

    {\Large Autorzy:\par}
    \vspace{0.5cm}
    {\large \authorone\par}
    \vspace{0.3cm}
    {\large \authortwo\par}

    \vfill

    {\large \projectdate\par}
\end{titlepage}

\tableofcontents
\newpage

% =====================================================================
\section{Wprowadzenie i opis problemu}
% =====================================================================
\textbf{Motywacja.} Celem niniejszego projektu jest opracowanie oraz implementacja algorytmu ewolucyjnego optymalizującego problem znany szerzej jako harmonogramowanie zadań w środowisku wieloprocesorowym z ograniczeniem pierwszeństwa (\textit{ang. Task Scheduling Problem with Precedence Constraints}). Jako że zadanie to należy do klasy problemów NP-trudnych~\cite{ullman1975}, naturalnym jest wykorzystanie metaheurystyk w celu jego rozwiązania.
\par\medskip\noindent
\textbf{Definicja formalna.} Problem zdefiniowany został jako czwórka $(T, p, l, M)$, gdzie:
\begin{itemize}[itemsep=3pt, topsep=3pt]
    \item $T = \{t_1, t_2, \dots, t_n\}$ jest zbiorem $n$ zadań do wykonania.
    \item $p : T \times T \to \{0, 1\}$ koduje relację pierwszeństwa: $p(t_i, t_j) = 1$ oznacza, że zadanie $t_i$ musi zostać zakończone przed rozpoczęciem zadania $t_j$, a wartość $0$ -- brak takiego więzu. Zakładamy, że graf skierowany indukowany przez~$p$ jest acykliczny (DAG); w~przeciwnym razie zbiór rozwiązań dopuszczalnych jest pusty.
    \item $l: T \to \mathbb{N}$ jest funkcją przypisującą każdemu zadaniu czas jego trwania.
    \item $M \in \mathbb{N}$ jest stałą liczbą grup pracowników, między które należy rozdzielić zadania. Grupa jest traktowana jako jedna jednostka, mogąca wykonać dowolne zadanie w jednym czasie nie łamiące relacji pierwszeństwa, dlatego będziemy się do niej odnosić jako procesor dla zgodności z literaturą i standardem.
\end{itemize}
\noindent
\textbf{Funkcja celu i ograniczenia.} Głównym wyzwaniem projektowym jest dokonanie takiego podziału zadań między $M$ procesorów, aby zachować spójność logiczną procesu. Funkcją celu jest minimalizacja czasu zakończenia wszystkich zadań, czyli tzw. \textit{makespan} ($C_{max}$). Optymalny harmonogram to taki, który przy pełnym wykorzystaniu dostępnych zasobów (grup) i zachowaniu wszystkich więzów technologicznych, kończy cały proces w najkrótszym możliwym czasie. W notacji Grahama problem zapisujemy jako $P\,|\,prec\,|\,C_{max}$.
\par\medskip\noindent
\textbf{Zakres dokumentu.} Niniejsza dokumentacja końcowa stanowi rozszerzenie dokumentacji wstępnej. Opisuje ona finalną metodę i jej implementację (rozdz.~\ref{sec:metoda}--\ref{sec:implementacja}), pełny plan i wyniki eksperymentów numerycznych (rozdz.~\ref{sec:eksperymenty}--\ref{sec:wyniki}) oraz wnioski wraz z weryfikacją postawionych hipotez badawczych (rozdz.~\ref{sec:dyskusja}--\ref{sec:wnioski}).

% =====================================================================
\section{Metoda rozwiązania}
\label{sec:metoda}
% =====================================================================
\textbf{Ogólna koncepcja.} W celu rozwiązania przedstawionego problemu zastosowano algorytm memetyczny~\cite{moscato1989}. Łączy on klasyczne dla algorytmów genetycznych globalne przeszukiwanie przestrzeni rozwiązań z optymalizacją lokalną opartą o analizę ścieżki krytycznej. Podejście to ma na celu zapobieganie powstawaniu wąskich gardeł w harmonogramach.

\subsection{Reprezentacja osobnika}
Do kodowania osobnika zastosowano reprezentację kolejności wykonywanych zadań. Chromosomem jest permutacja o długości $|T|$, zawierająca wszystkie zadania w kolejności priorytetu, z jakim będą umieszczane w harmonogramie. Reprezentacja ta nie zawiera bezpośredniego przydziału do grup -- przydział wynika z działania dekodera. Każda permutacja zgodna z porządkiem topologicznym koduje rozwiązanie dopuszczalne.

\subsection{Dekoder (\textit{list scheduling})}
W celu zamiany chromosomu na faktyczny harmonogram oraz wyliczenia czasu jego trwania zastosowano zachłanny dekoder typu \textit{list scheduling}~\cite{graham1969}. Ponieważ procesory są nierozróżnialne, dekoder po kolei przydziela wolnym jednostkom kolejne gotowe zadania (zgodnie z priorytetem z chromosomu), respektując relację pierwszeństwa~$p$. Wyliczona przez dekoder wartość $C_{max}$ stanowi wartość funkcji przystosowania. Dekoder udostępnia również pełny harmonogram (czasy rozpoczęcia/zakończenia oraz przydział do procesora każdego zadania), co jest wykorzystywane przez moduł optymalizacji lokalnej.

\subsection{Naprawa dopuszczalności (\textit{ready-list repair})}
Operatory krzyżowania i mutacji mogą wygenerować permutacje niezgodne z porządkiem topologicznym. Aby zapewnić, że każdy osobnik koduje rozwiązanie dopuszczalne, zastosowano mechanizm naprawy oparty o listę zadań gotowych: chromosom jest przebudowywany do najbliższego porządku topologicznego z zachowaniem względnej kolejności z oryginalnej permutacji. Dzięki temu cała populacja pozostaje dopuszczalna i nie jest konieczne karanie rozwiązań niedopuszczalnych w funkcji celu.

\subsection{Operatory ewolucyjne}
\begin{itemize}[itemsep=3pt, topsep=3pt]
    \item \textbf{Selekcja:} selekcja turniejowa~\cite{miller1995}, mająca zapobiegać przedwczesnej dominacji pojedynczego osobnika.
    \item \textbf{Krzyżowanie:} z uwagi na permutacyjny charakter reprezentacji zastosowano operatory unikające dublowania zadań w chromosomie. Porównano krzyżowanie porządkowe (Order Crossover, OX)~\cite{davis1985}, zachowujące względną kolejność zadań, z krzyżowaniem z częściowym mapowaniem (PMX)~\cite{goldberg1985}, opartym o mieszanie pozycyjne.
    \item \textbf{Mutacja:} operator zamiany (\textit{swap}) -- losowa zamiana miejscami dwóch zadań w chromosomie.
\end{itemize}

\subsection{Optymalizacja lokalna (CPM)}
Kluczowym elementem metody memetycznej jest moduł optymalizacji lokalnej oparty o analizę ścieżki krytycznej (\textit{ang. Critical Path Method})~\cite{kelley1959}. W odróżnieniu od klasycznej CPM liczonej wyłącznie na grafie zależności, zastosowano \textbf{analizę świadomą harmonogramu}: ścieżka krytyczna wyznaczana jest na faktycznie zdekodowanym harmonogramie, śledząc na każdym kroku poprzednika wynikającego z relacji pierwszeństwa \emph{lub} z zajętości procesora. Dzięki temu analiza wykrywa również wąskie gardła zasobowe i różni się między osobnikami. Zadania leżące na ścieżce krytycznej otrzymują wyższy priorytet (operacja typu \textit{pull-left}); jeśli taka modyfikacja poprawia $C_{max}$, zgodnie z ewolucją Lamarckowską~\cite{whitley1994} zastępuje ona poprzednika w populacji.

\subsection{Konfiguracje algorytmu}
Z połączenia dwóch operatorów krzyżowania (OX/PMX) oraz obecności bądź braku modułu CPM powstają cztery konfiguracje porównywane w eksperymentach: \texttt{ox}, \texttt{pmx} (czysty AE) oraz \texttt{ox\_cpm}, \texttt{pmx\_cpm} (pełny algorytm memetyczny).

% =====================================================================
\section{Implementacja}
\label{sec:implementacja}
% =====================================================================
\textbf{Technologie.} Projekt zrealizowano w języku Python (wersja 3.10+), wykorzystując bibliotekę \textit{DEAP}~\cite{deap2012} jako szkielet algorytmu ewolucyjnego. Rozwiązania optymalne dla małych instancji wyznaczono solverem MILP (CBC) z pakietu \textit{PuLP}~\cite{mitchell2011}. Pozostałe komponenty oparto o standardowe biblioteki naukowe: \textit{NumPy}~\cite{harris2020}, \textit{Matplotlib}~\cite{hunter2007} oraz \textit{SciPy}~\cite{virtanen2020} (testy statystyczne).
\par\medskip\noindent
\textbf{Struktura repozytorium.} Najważniejsze moduły:
\begin{itemize}[itemsep=2pt, topsep=3pt]
    \item \texttt{src/core/} -- model dziedzinowy, dekoder \textit{list scheduling}, analiza CPM świadoma harmonogramu;
    \item \texttt{src/data/} -- parser instancji STG (Kasahara Lab) oraz generator grafów syntetycznych;
    \item \texttt{src/solvers/evolutionary/} -- algorytm memetyczny (operatory, strategie, lokalne przeszukiwanie CPM);
    \item \texttt{src/solvers/heuristics.py} -- heurystyki listowe HLFET i MCP;
    \item \texttt{src/solvers/milp.py} -- dokładny model MILP (czasowo-indeksowany $P\,|\,prec\,|\,C_{max}$);
    \item \texttt{src/experiments/} -- runner kampanii eksperymentalnej, analiza statystyczna i wykresy.
\end{itemize}
\par\medskip\noindent
\textbf{Powtarzalność.} Ustalono ziarna generatora liczb losowych, a środowisko (wersje bibliotek) udokumentowano w pliku \texttt{requirements.txt}. Pełną kampanię odtwarza się poleceniami:
\begin{lstlisting}[language=bash]
python -m src.experiments.run_experiments   # surowe wyniki + krzywe zbieznosci
python -m src.experiments.analyze           # podsumowanie + testy Wilcoxona (H1/H2)
python -m src.experiments.plots             # wykresy pudelkowe i krzywe zbieznosci
\end{lstlisting}

\todo{ewentualnie dodać krótki opis najciekawszych decyzji implementacyjnych / problemów napotkanych w trakcie (np. naprawa CPM świadomej harmonogramu, kara za niedopuszczalność jako martwy kod).}

% =====================================================================
\section{Plan eksperymentów numerycznych}
\label{sec:eksperymenty}
% =====================================================================
\textbf{Cel i hipotezy badawcze.} Celem eksperymentów jest ocena jakości rozwiązań zaproponowanego algorytmu oraz weryfikacja dwóch hipotez badawczych:
\begin{itemize}[itemsep=1pt, topsep=2pt, leftmargin=3em]
    \item[(H1)] lokalna optymalizacja CPM istotnie obniża medianę $C_{max}$ względem czystego AE;
    \item[(H2)] krzyżowanie OX istotnie dominuje nad PMX na grafach rzadkich ($|E|/|T|^2 < 0{,}1$).
\end{itemize}
\par\medskip\noindent
\textbf{Zestawy testowe.} Jako instancje testowe posłużyły grafy ze zbioru \textit{Standard Task Graph Set} (Kasahara Lab)~\cite{tobita2002} oraz grafy generowane syntetycznie z kontrolowaną liczbą zadań i gęstością. Dla małych instancji ($|T| \leq 30$) wyznaczono rozwiązanie optymalne solverem MILP w celu raportowania luki optymalizacyjnej; dla większych instancji odniesieniem jest dolne ograniczenie $LB = \max(\lceil \text{całkowita praca}/M \rceil,\ \text{ścieżka krytyczna})$.
\par\medskip\noindent
\textbf{Protokół statystyczny.} Dla każdej konfiguracji wykonano co najmniej 30 niezależnych uruchomień (10 dla dwóch największych instancji ze względów obliczeniowych). Miarą jakości jest najlepszy $C_{max}$ na uruchomienie; tabele raportują medianę. Istotność różnic weryfikowano sparowanym testem kolejności rangowanej Wilcoxona~\cite{wilcoxon1945} (te same ziarna), z korekcją Holma w obrębie każdej rodziny hipotez, $\alpha = 0{,}05$.
\par\medskip\noindent
\textbf{Baseline'y i ablacje.} Podejście zestawiono z klasycznymi heurystykami listowymi HLFET~\cite{adam1974} oraz MCP~\cite{wu1990}, a także z czystym AE pozbawionym komponentu lokalnego.

\paragraph{Zastrzeżenie dot. budżetu (H1).} Konfiguracje memetyczne wykonują dodatkowe ewaluacje funkcji celu wewnątrz lokalnego przeszukiwania CPM, dlatego porównanie prowadzone jest przy równej liczbie \emph{pokoleń}, a nie przy równym budżecie ewaluacji. Przewaga CPM jest zatem częściowo skutkiem większego efektywnego nakładu obliczeniowego na pokolenie. \todo{ewentualne porównanie przy równym budżecie ewaluacji -- jeśli zostanie wykonane.}

\paragraph{Hiperparametry.} \todo{uzupełnić finalne wartości: rozmiar populacji, liczba pokoleń, $p_c$, $p_m$, rozmiar turnieju. Opisać ewentualne ablacje hiperparametrów lub zaznaczyć, że nie były przeprowadzone.}

% =====================================================================
\section{Wyniki}
\label{sec:wyniki}
% =====================================================================
\todo{Cały rozdział do finalizacji po ostatecznym przebiegu kampanii. Liczby poniżej pochodzą z bieżącego przebiegu i należy je potwierdzić/odświeżyć. Źródła: \texttt{results/campaign/summary.csv}, \texttt{h1\_wilcoxon.csv}, \texttt{h2\_wilcoxon.csv} oraz \texttt{RESULTS.md}.}

\subsection{Charakterystyka instancji}
\todo{Wstawić finalną tabelę instancji: rozmiar, liczba krawędzi, gęstość, $M$, HLFET, MCP, optimum MILP (jeśli dostępne), dolne ograniczenie LB. Wzorzec poniżej.}

\begin{table}[h]
\centering
\small
\begin{tabular}{lrrrrr}
\toprule
instancja & $|T|$ & $|E|$ & gęstość & $M$ & ref. ($C_{max}$) \\
\midrule
small\_15 & 17 & 31 & 0{,}107 & 2 & 52 (MILP) \\
small\_22 & 24 & 48 & 0{,}083 & 3 & 44 (MILP) \\
\todo{\dots} & & & & & \\
\bottomrule
\end{tabular}
\caption{Charakterystyka instancji testowych. \todo{uzupełnić pełny zestaw}}
\label{tab:instances}
\end{table}

\subsection{Jakość rozwiązań -- mediana $C_{max}$ wg konfiguracji}
\todo{Wstawić finalną tabelę mediany $C_{max}$ dla czterech konfiguracji (ox / pmx / ox\_cpm / pmx\_cpm) wraz z luką do odniesienia. Najlepszą konfigurację na instancję zaznaczyć pogrubieniem.}

\begin{table}[h]
\centering
\small
\begin{tabular}{lrrrrr}
\toprule
instancja & ref. & ox & pmx & ox\_cpm & pmx\_cpm \\
\midrule
\todo{\dots} & & & & & \\
\bottomrule
\end{tabular}
\caption{Mediana $C_{max}$ według konfiguracji. \todo{uzupełnić}}
\label{tab:cmax}
\end{table}

\subsection{Weryfikacja H1 -- wpływ optymalizacji CPM}
\todo{Wstawić finalną tabelę testów Wilcoxona CPM-on vs CPM-off (per krzyżowanie, per instancja), z $\Delta$ mediany i skorygowaną $p$. Skomentować werdykt.}

\subsection{Weryfikacja H2 -- OX vs PMX na grafach rzadkich}
\todo{Wstawić finalną tabelę testów Wilcoxona OX vs PMX (czysty AE), z podziałem na grafy rzadkie/gęste. Skomentować werdykt.}

\subsection{Wykresy}
\todo{Wstawić wykresy pudełkowe $C_{max}$ oraz krzywe zbieżności dla reprezentatywnych instancji z katalogu \texttt{results/campaign/plots/}. Przykładowy szkielet poniżej.}

% \begin{figure}[h]
%   \centering
%   \includegraphics[width=0.49\textwidth]{../../results/campaign/plots/random_300_boxplot.png}
%   \includegraphics[width=0.49\textwidth]{../../results/campaign/plots/random_300_convergence.png}
%   \caption{random_300: wykres pudełkowy oraz krzywa zbieżności.}
%   \label{fig:random300}
% \end{figure}

% =====================================================================
\section{Dyskusja}
\label{sec:dyskusja}
% =====================================================================
\todo{Do finalizacji. Punkty do omówienia (na podstawie bieżących wyników -- potwierdzić):}
\begin{itemize}[itemsep=2pt, topsep=3pt]
    \item \textbf{H1} -- wstępnie potwierdzona: CPM istotnie obniża $C_{max}$ wszędzie tam, gdzie istnieje zapas (z zastrzeżeniem dot. budżetu pokoleń); efekt rośnie z trudnością instancji. \todo{potwierdzić finalnymi liczbami.}
    \item \textbf{H2} -- wstępnie odrzucona, zachodzi relacja przeciwna: PMX istotnie dominuje nad OX na grafach rzadkich. Prawdopodobna przyczyna: naprawa listy gotowych już wymusza pierwszeństwo, więc przewaga OX w zachowaniu kolejności jest stłumiona, podczas gdy pozycyjne mieszanie PMX skuteczniej eksploruje przestrzeń priorytetów. \todo{potwierdzić finalnymi liczbami.}
    \item Algorytm memetyczny pokonuje heurystyki listowe HLFET/MCP na każdej instancji; pozostała luka do dolnego ograniczenia. \todo{podać finalne wartości luk.}
    \item Ograniczenia metody i zagrożenia trafności (brak elityzmu, porównanie przy równej liczbie pokoleń, dobór instancji). \todo{rozwinąć.}
\end{itemize}

% =====================================================================
\section{Wnioski}
\label{sec:wnioski}
% =====================================================================
\todo{Do finalizacji po zatwierdzeniu wyników. Podsumować: najlepsza konfiguracja (wstępnie \texttt{pmx\_cpm}), werdykty H1/H2, przewaga nad baseline'ami, kierunki dalszych prac (ablacje hiperparametrów, elityzm, porównanie przy równym budżecie ewaluacji).}

% =====================================================================
\begin{thebibliography}{99}

\bibitem{moscato1989}
P.~Moscato,
\textit{On Evolution, Search, Optimization, Genetic Algorithms and Martial Arts: Towards Memetic Algorithms},
Caltech Concurrent Computation Program, Report C3P 826, 1989.

\bibitem{graham1969}
R.~L. Graham,
\textit{Bounds on multiprocessing timing anomalies},
SIAM Journal on Applied Mathematics, vol.~17, no.~2, pp.~416--429, 1969.

\bibitem{davis1985}
L.~Davis,
\textit{Applying adaptive algorithms to epistatic domains},
w: Proc. of the 9th International Joint Conference on Artificial Intelligence (IJCAI), pp.~162--164, 1985.

\bibitem{goldberg1985}
D.~E. Goldberg, R.~Lingle,
\textit{Alleles, loci, and the traveling salesman problem},
w: Proc. of the 1st International Conference on Genetic Algorithms, pp.~154--159, 1985.

\bibitem{kelley1959}
J.~E. Kelley, M.~R. Walker,
\textit{Critical-path planning and scheduling},
w: Proc. of the Eastern Joint Computer Conference, pp.~160--173, 1959.

\bibitem{tobita2002}
T.~Tobita, H.~Kasahara,
\textit{A standard task graph set for fair evaluation of multiprocessor scheduling algorithms},
Journal of Scheduling, vol.~5, no.~5, pp.~379--394, 2002.
\url{https://www.kasahara.cs.waseda.ac.jp/schedule/}

\bibitem{adam1974}
T.~L. Adam, K.~M. Chandy, J.~R. Dickson,
\textit{A comparison of list schedules for parallel processing systems},
Communications of the ACM, vol.~17, no.~12, pp.~685--690, 1974.

\bibitem{wu1990}
M.-Y. Wu, D.~D. Gajski,
\textit{Hypertool: A programming aid for message-passing systems},
IEEE Transactions on Parallel and Distributed Systems, vol.~1, no.~3, pp.~330--343, 1990.

\bibitem{wilcoxon1945}
F.~Wilcoxon,
\textit{Individual comparisons by ranking methods},
Biometrics Bulletin, vol.~1, no.~6, pp.~80--83, 1945.

\bibitem{deap2012}
F.-A. Fortin, F.-M. De~Rainville, M.-A. Gardner, M.~Parizeau, C.~Gagné,
\textit{DEAP: Evolutionary algorithms made easy},
Journal of Machine Learning Research, vol.~13, pp.~2171--2175, 2012.

\bibitem{ullman1975}
J.~D. Ullman,
\textit{NP-complete scheduling problems},
Journal of Computer and System Sciences, vol.~10, no.~3, pp.~384--393, 1975.

\bibitem{miller1995}
B.~L. Miller, D.~E. Goldberg,
\textit{Genetic algorithms, tournament selection, and the effects of noise},
Complex Systems, vol.~9, no.~3, pp.~193--212, 1995.

\bibitem{whitley1994}
D.~Whitley, V.~S. Gordon, K.~Mathias,
\textit{Lamarckian evolution, the Baldwin effect and function optimization},
w: Parallel Problem Solving from Nature (PPSN III), LNCS 866, Springer, pp.~6--15, 1994.

\bibitem{harris2020}
C.~R. Harris \textit{et al.},
\textit{Array programming with NumPy},
Nature, vol.~585, no.~7825, pp.~357--362, 2020.

\bibitem{hunter2007}
J.~D. Hunter,
\textit{Matplotlib: A 2D graphics environment},
Computing in Science \& Engineering, vol.~9, no.~3, pp.~90--95, 2007.

\bibitem{virtanen2020}
P.~Virtanen \textit{et al.},
\textit{SciPy 1.0: Fundamental algorithms for scientific computing in Python},
Nature Methods, vol.~17, pp.~261--272, 2020.

\bibitem{mitchell2011}
S.~Mitchell, M.~O'Sullivan, I.~Dunning,
\textit{PuLP: A linear programming toolkit for Python},
Technical report, University of Auckland, 2011.
\url{https://github.com/coin-or/pulp}

\end{thebibliography}

\end{document}
